\section{Keeping the same field space in each phase}
\label{ex-phase-fields}

A field label may disappear when one passes to a component of a mixture.
For the selected phase to represent the intended fields, one must rule out
additional null labels. The following argument explains the role of the
finite detectors and the uniform moment estimates used for this purpose.

\begin{lemma}[Moment bounds at every character]\label{ex-character-moments}
Let \(A\in\mathcal A\), let \(t\ge0\), and suppose
\[
 \langle\Omega,(A^*A)^n\Omega\rangle\le t^{2n}
 \qquad(n=0,1,2,\ldots).
\]
Then \(|\chi(A)|\le t\) for every \(\chi\in\Xi\).
\end{lemma}
\begin{proof}
The representing measure gives
\(\int|\chi(A)|^{2n}\,d\mu\le t^{2n}\).
If \(|\chi_0(A)|>t\), choose \(p\) strictly between these numbers.
The nonempty open set \(O=\{\chi:|\chi(A)|>p\}\) has positive measure
\(c\) by full support. Hence
\(cp^{2n}\le\int|\chi(A)|^{2n}\,d\mu\le t^{2n}\).
Since \(p>0\), division gives \(c\le(t/p)^{2n}\to0\), a
contradiction. This also covers \(t=0\), when the right side already
vanishes for positive \(n\).
\end{proof}

Here is the elementary two-point calculation to which these moment
estimates are applied. At a fixed scale \(k\), let \(a,b\) denote two
insertions at distinct points and let \(j_a,j_b,\kappa\in\mathbb C\)
be the prescribed comparison one-point values and connected two-point
value. Form the elements
\begin{align*}
 D_a&=M(a)-j_aI,\qquad D_b=M(b)-j_bI,\\
 D_{ab}&=M(ab)-j_aM(b)-j_bM(a)+(j_aj_b-\kappa)I.
\end{align*}
Assume the preceding lemma's moment bounds hold for each of these three
elements with the same number \(\epsilon_k\ge0\). It follows for every
character that
\begin{align}
 |S_\chi(a)-j_a|&\le\epsilon_k,\label{ex-detector-one}\\
 |S_\chi(ab)-(\kappa+j_aj_b)|
 &\le\epsilon_k(1+|j_a|+|j_b|).\label{ex-detector-two}
\end{align}
Indeed, evaluating \(D_{ab}\) and rearranging gives
\[
 S_\chi(ab)-(\kappa+j_aj_b)
 =\chi(D_{ab})+j_a(S_\chi(b)-j_b)+j_b(S_\chi(a)-j_a),
\]
so the triangle inequality proves \eqref{ex-detector-two}.
If \(\epsilon_k\to0\) and the prescribed one-point values converge,
these errors tend to zero in every phase, with a bound independent of
the character. This conclusion requires moments of the full corrected
operators \(D_{ab}\), not only pointwise estimates for the unselected
two-point correlations.

\begin{proposition}[Finite detectors prevent an additional null space]
\label{ex-no-new-null}
Let \(E\) be a complex vector space of raw labels, and let
\(N\subset N_\chi\) be subspaces. Write
\(\pi:E\to E/N=V\) and \(\pi_\chi:E\to E/N_\chi=V_\chi\).
Let \(W\) be a finite-dimensional complex normed space. Suppose there
are linear maps \(J_k:W\to E\), linear tests
\(\ell_{k,i}:V_\chi\to\mathbb C\), \(1\le i\le r\), and an
injective linear map \(F:W\to\mathbb C^r\) such that:
\begin{enumerate}
\item \(\pi J_k\) is surjective for all sufficiently large \(k\);
\item for every vector of a fixed basis of \(W\),
\[
 F_k w:=\bigl(\ell_{k,i}(\pi_\chi J_kw)\bigr)_{i=1}^r
 \longrightarrow Fw.
\]
\end{enumerate}
Then \(N_\chi=N\).
\end{proposition}
\begin{proof}
Convergence on a finite basis implies convergence in operator norm.
To see this, let \(e_1,\ldots,e_s\) be that basis with continuous
coordinate functionals \(e_j^*\). Then
\[
 \|(F_k-F)w\|
 \le\sum_{j=1}^s|e_j^*(w)|\,\|(F_k-F)e_j\|
 \le\|w\|\sum_{j=1}^s\|e_j^*\|\,\|(F_k-F)e_j\|.
\]
The finite sum tends to zero. If \(W\ne\{0\}\), compactness of its
unit sphere and injectivity of \(F\) give
\(c=\min_{\|w\|=1}\|Fw\|>0\). Choose \(k\) late enough that
\(\|F_k-F\|<c/2\) and \(\pi J_k\) is onto. Then
\(\|F_kw\|\ge(c/2)\|w\|\), so \(F_k\) is injective.
If \(W=\{0\}\), surjectivity already gives \(V=0\), and the result
is immediate.

Now take \(e\in N_\chi\). By surjectivity choose \(w\in W\) with
\(\pi J_kw=\pi e\). Thus \(J_kw-e\in N\subset N_\chi\), and since
\(e\in N_\chi\), also \(J_kw\in N_\chi\). All tests vanish on
\(\pi_\chi J_kw=0\); hence \(F_kw=0\). Injectivity forces \(w=0\),
so \(\pi e=\pi J_k0=0\), or \(e\in N\). This proves
\(N_\chi\subset N\); the reverse inclusion was assumed.
\end{proof}

In the application at label degree \(d\), \(W\) is the comparison
field space, \(J_k\) consists of its corrected labels at the chosen
scales, and the tests \(\ell_{k,i}\) are one- or two-point correlations
whose exterior labels are the corrected labels at comparison index \(k\);
they depend on \(k\), as the proposition allows. The test points are
\(x_{k,i}=\sigma_ky_{k,i}\), where \(\sigma_k\to0\) is the comparison scale
and the rescaled configurations \(y_{k,i}\) converge to configurations
\(y_i\) of distinct points; the physical configurations \(x_{k,i}\) do not
converge to a configuration of distinct points.
Estimates~\eqref{ex-detector-one}--\eqref{ex-detector-two} compare the
phase correlations at \(x_{k,i}\) with the free values at \(y_{k,i}\), with
errors \(\epsilon_k\to0\) uniform in the character. Continuity of the free
comparison table at the distinct configurations \(y_i\) then gives
convergence of \(F_k\) to \(F\) on a basis, where \(F\) is given by the free
tests at \(y_i\). The comparison limit \(k\to\infty\) is taken after the
mesh and period limits that define the table. The comparison tests separate \(W\), so
their joint map \(F\) is injective. The first hypothesis of the
proposition is the eventual spanning statement for the corrected labels
in the unselected quotient. The same-label construction of
Section~\ref{sec:actual-rank-purchases}, followed by the exact-kernel
argument in Section~\ref{sec:finite-rank-kernel-transport}, supplies
this spanning statement. Section~\ref{sec:actual-phase-detector-moments}
supplies the corrected-operator moment bounds on those trials, and
Theorem~\ref{thm:ex-lattice-all-grade-free-detector} supplies the
separating free tests. These applications retain the same corrected
lifts and all separated correlations when passing to the smooth table.
Lemma~\ref{ex-phase-original-null} supplies $N\subseteq N_\chi$.
The proposition then proves equality in every grade for every character.
Thus the finite-dimensional argument uses the actual comparison and
moment constructions; finite dimensionality alone does not supply them.

\subsection{Field inclusions from simultaneous tested representatives}
\label{subsec:phase-field-inclusions}

Equality of the null spaces at each grade does not construct a map between
those grades. We now give that construction from a separate realization
input, on one fixed table. No change of phase or of the extraction defining
the table is made in this argument.

Let $A_d$ be a finite-dimensional complex raw label space for each
$d\in\mathbb N$. A raw label includes its grade. Fix a table $S$ whose
correlations
\[
 S(x;(d_1,a_1),\ldots,(d_m,a_m)),\qquad
 x\in\operatorname{Conf}_m(\mathbb R^4),\quad a_j\in A_{d_j},
\]
are complex multilinear in the labels, smooth in $x$, and symmetric under
simultaneous permutation of labels and positions. Here
$\operatorname{Conf}_m(\mathbb R^4)$ is the open set of pairwise distinct
positions. Spectator tuples may have arbitrary finite length and arbitrary
grades; no upper grade or arity bound is imposed on them.

An exterior $\xi=(x;\boldsymbol b)$ contains a tuple $\boldsymbol b$
of $n$ raw spectators and a configuration of $n+1$ distinct positions,
$x\in\operatorname{Conf}_{n+1}(\mathbb R^4)$. Allow $n=0$. Let $\mathcal X$ be the set of all such exteriors and define the
linear response map
\[
 \mathcal R_d:A_d\longrightarrow\mathbb C^{\mathcal X},\qquad
 (\mathcal R_da)(x;\boldsymbol b)=S(x;(d,a),\boldsymbol b).
\]
Set
\[
 N_d=\ker\mathcal R_d,\qquad E_d=A_d/N_d,
 \qquad \pi_d:A_d\longrightarrow E_d.
\]
Thus nullity means vanishing against every exterior, not just a finite
detector family. The spaces $E_d$ are finite-dimensional. The response map
descends to an injective linear map
$J_d:E_d\to\mathbb C^{\mathcal X}$: if $J_d\pi_da=0$, then
$a\in N_d$, so $\pi_da=0$.

Choose a basis $(a_{d,j})_{j\in B_d}$ of each $A_d$. The incoming
\emph{simultaneous tested realization} is the following statement. For
every $d\le e$ and every $j\in B_d$, there is one $b_{de,j}\in A_e$
such that, for every spectator tuple $\boldsymbol b$ of length $n$ and
every $\varphi\in C_c^\infty(\operatorname{Conf}_{n+1}(\mathbb R^4))$,
\begin{equation}\label{eq:phase-inclusion-tested-input}
 \int\varphi(x)S(x;(e,b_{de,j}),\boldsymbol b)\,dx
 =\int\varphi(x)S(x;(d,a_{d,j}),\boldsymbol b)\,dx.
\end{equation}
The order of these quantifiers matters: $b_{de,j}$ is chosen before the
spectators and the test. A representative depending on either would not
supply a linear map of label spaces. This input asks only for raw
representatives of insertion functionals; it assumes neither quotient
inclusions nor their injectivity or coherence.

\begin{proposition}[Conditional construction of the field inclusions]
\label{prop:phase-field-inclusions}
For the fixed smooth symmetric table above, the simultaneous tested
realization supplies unique linear maps
$\iota_{de}:E_d\to E_e$ for $d\le e$ preserving every exterior response.
They are injective and satisfy
\[
 \iota_{dd}=1,\qquad
 \iota_{ef}\iota_{de}=\iota_{df}\quad(d\le e\le f).
\]
Correlations descend to the quotient labels in every slot. For any fixed
configuration of distinct points, any slot, and arbitrary quotient
spectators, replacing a label $v\in E_d$ in that slot by
$\iota_{de}v\in E_e$ leaves the correlation unchanged.
\end{proposition}
\begin{proof}
First define a raw linear map by basis expansion:
\[
 K_{de}\left(\sum_{j\in B_d}c_ja_{d,j}\right)
       =\sum_{j\in B_d}c_jb_{de,j}.
\]
For each fixed spectator tuple and test, multilinearity and the finite sum
in this definition extend \eqref{eq:phase-inclusion-tested-input} to every
$a\in A_d$, with $b_{de,j}$ replaced by $K_{de}a$. The same map works for
all tests and all spectators because its basis images were fixed first.

We spell out why this tested equality is pointwise. For a fixed $a$ and
spectator tuple, let
\[
 h(x)=S(x;(e,K_{de}a),\boldsymbol b)
                   -S(x;(d,a),\boldsymbol b).
\]
It is continuous on the open configuration space and has zero integral
against every smooth compactly supported test there. If $h(x_0)\ne0$,
choose $c\in\mathbb C$ with $|c|=1$ and $c h(x_0)=|h(x_0)|$.
There is a ball about $x_0$, with closure in the configuration space,
on which $\operatorname{Re}(c h)>|h(x_0)|/2$. Choose a nonnegative smooth
function $\psi$, supported in that ball, with positive integral. Such a
function is obtained by translating and rescaling
$\exp(-1/(1-|x|^2))$ on $|x|<1$, extended by zero outside. Then
\[
 0=\operatorname{Re}\left(c\int\psi h\right)
   >\frac{|h(x_0)|}{2}\int\psi>0,
\]
a contradiction. Thus $h=0$ on the configuration space. In particular,
\begin{equation}\label{eq:phase-inclusion-raw-response}
                 \mathcal R_eK_{de}=\mathcal R_d.
\end{equation}
No equality at coincident points is used or asserted.

If $a\in N_d$, this identity gives $K_{de}a\in N_e$. Consequently
\[
               \iota_{de}(\pi_da)=\pi_eK_{de}a
\]
is well defined and linear. If different simultaneous representatives, or
a different basis, give a raw map $K'_{de}$ with the same tested property,
then \eqref{eq:phase-inclusion-raw-response} holds for both maps. Their
difference takes values in $N_e$, so they induce the same quotient map.
For every $v\in E_d$ we have
\begin{equation}\label{eq:phase-inclusion-faithful-response}
                     J_e\iota_{de}v=J_dv.
\end{equation}
Since $J_d$ is injective, this also proves injectivity of $\iota_{de}$.
Since $J_e$ is injective, it proves uniqueness of a response-preserving
map $E_d\to E_e$.

For completeness, null labels vanish in every slot. Permute the chosen
slot and its position to the first position; symmetry expresses the
correlation as a head insertion against another raw exterior. If the
chosen label belongs to its $N_d$, that head correlation is zero by
definition. Thus changing any raw representative by a null label changes
no correlation. Changing the finitely many representatives one at a time
proves well-defined multilinear correlations on quotient tuples. Smoothness
and permutation symmetry pass to these correlations by taking fixed raw
representatives. The pointwise identity already proved at the head,
followed by the same permutation, gives explicitly
\begin{equation}\label{eq:phase-inclusion-slot}
 S(x;v_1,\ldots,\iota_{de}v_i,\ldots,v_m)
       =S(x;v_1,\ldots,v_i,\ldots,v_m),
 \qquad v_i\in E_d,
\end{equation}
where the other labels have any fixed grades and the changed slot is
assigned grade $e$ on the left. Repeated application gives the identity
when any finite collection of slots is raised, even to different grades.
Integration against a compactly supported test preserves the identity.

Finally $J_d\iota_{dd}=J_d$ implies $\iota_{dd}=1$. If $d\le e\le f$,
then
\[
 J_f\iota_{ef}\iota_{de}=J_e\iota_{de}=J_d
                         =J_f\iota_{df}.
\]
Injectivity of $J_f$ proves composition. These identities are conclusions
on the quotients; the chosen raw maps need not compose or equal the raw
identity when $d=e$.
\end{proof}

\begin{proposition}[The common union and its universal property]
\label{prop:phase-field-union}
Under the hypotheses of Proposition~\ref{prop:phase-field-inclusions},
the spaces $E_d$ have an algebraic direct limit $E$, with injective maps
$j_d:E_d\to E$ and $j_e\iota_{de}=j_d$. Every element of $E$ comes from
one grade, and every finite collection comes from one common grade. The
table correlations are well defined on finite tuples from $E$.
If $W$ is any complex vector space and linear maps $F_d:E_d\to W$ satisfy
$F_e\iota_{de}=F_d$ for all $d\le e$, there is a unique linear
$F:E\to W$ such that $Fj_d=F_d$ for every $d$.
\end{proposition}
\begin{proof}
Realize the union in the single response space:
\[
               E=\bigcup_{d\ge0}J_d(E_d)
                    \ \subseteq\mathbb C^{\mathcal X}.
\]
Equation~\eqref{eq:phase-inclusion-faithful-response} makes these ranges
increasing. If $u=J_dv$ and $w=J_ez$, put $D=\max(d,e)$. For any scalars
$a,b\in\mathbb C$,
\[
 au+bw=J_D\bigl(a\iota_{dD}v+b\iota_{eD}z\bigr)\in E.
\]
The zero vector belongs to every range, so $E$ is a vector subspace.
Take $j_d$ to be $J_d$ with codomain restricted to $E$. These maps are
injective and satisfy the stated compatibility. If $u_1,\ldots,u_m\in E$,
choose $u_i=j_{d_i}v_i$ and let $D=\max_i d_i$. Then
$u_i=j_D\iota_{d_iD}v_i$ for every $i$. For the empty collection choose
$D=0$. Surjectivity of $\pi_D$ also gives simultaneous raw representatives
in $A_D$ for any such finite collection. No bound uniform over all elements
of $E$ is implied.

To evaluate a correlation of $u_1,\ldots,u_m$, choose these representatives
at a common grade and use the quotient correlation constructed above.
Suppose a second choice uses grade $D'$. At $L=\max(D,D')$, the two
representatives of each $u_i$ have the same image under the injective map
$j_L$ and hence are equal in $E_L$. Raising all slots to $L$ leaves the
correlation unchanged by \eqref{eq:phase-inclusion-slot}. Thus the value
is independent of both the common grade and all representatives. The same
argument treats initially different grades slot by slot. Multilinearity,
symmetry and smoothness follow by evaluating each finite tuple at one
fixed common grade; the zero-point correlation is the original vacuum
value and requires no label representative.

Define $F(j_dv)=F_dv$. If $j_dv=j_ez$, then at $D=\max(d,e)$,
\[
 j_D\iota_{dD}v=j_D\iota_{eD}z,
 \qquad\hbox{hence}\qquad \iota_{dD}v=\iota_{eD}z.
\]
Compatibility of the $F_d$ proves $F_dv=F_ez$, so $F$ is well defined.
Putting two arguments in one grade proves additivity and scalar
homogeneity from those of $F_D$. Its restriction to each grade is $F_d$.
Any other such map agrees on every $j_d(E_d)$, whose union is $E$, and
therefore equals $F$. This is the algebraic direct-limit property of
complex vector spaces; no topology or multiplication on the union is
being asserted.
\end{proof}

These conclusions apply to a selected phase only with the realization
receipt for that same selected table. More generally, suppose a second
smooth symmetric table $T$ on the same raw spaces has
$N_d^S\subseteq N_d^T$ at every grade, and the same raw maps $K_{de}$
satisfy \eqref{eq:phase-inclusion-tested-input} for both tables. The natural
surjections $q_d:E_d^S\to E_d^T$, $q_d\pi_d^Sa=\pi_d^Ta$, then obey
\[
 q_e\iota_{de}^S\pi_d^Sa=\pi_e^TK_{de}a
             =\iota_{de}^Tq_d\pi_d^Sa.
\]
Surjectivity of $\pi_d^S$ proves compatibility on all quotient labels.
The universal property gives the induced map between the two unions.
Correlations of $T$ on the $S$-quotients are well defined by null-space
containment, and their inclusion compatibility follows from the same
pointwise tested identity for $T$. If the null spaces are equal, each
$q_d$ is an isomorphism, and the induced union map is an isomorphism:
its inverse is obtained from the compatible inverses $q_d^{-1}$.
The numerical correlations of the two tables need not be equal.
Without the tested realization for $T$ on these same raw maps, grade-wise
null-space equality alone supplies none of these cross-grade correlation
identities. Derivative maps and their slot identities, and compatibility
with calibrated actions or involutions, require their own receipts and
are not constructed in this subsection.

\subsection{Attaching one table and passing its inclusions to every character}
\label{subsec:phase-inclusion-inheritance}

\paragraph{The same-table realization input.}
\label{par:phase-inclusion-attachment}
There is a conditional supplier for
\eqref{eq:phase-inclusion-tested-input} in the preceding constructions.
Fix one finite-volume family, its raw reference spaces $A_d$, its terminal
calibrations, and one retained extraction. Let $S$ be the smooth symmetric
table representing the ordered smeared limits of that same family. All the
hypotheses of the insertion comparison in
Section~\ref{sec:actual-insertion-derivative-relay}, the projection and
all-test kernel construction in
Section~\ref{sec:finite-rank-kernel-transport}, and the normalized test
argument in Section~\ref{sec:ope-filtration} are retained on these data.
In particular, this retains their local service and equation inputs,
comparison and rank purchases, inverse-calibration bounds, and convergence
of the smeared correlations to this $S$. The ordered limit permits the
mesh threshold to depend on the selected period; no mesh threshold uniform in the period
or diagonal limit is substituted.

Here is the precise application. Fix $d\le e$ and one basis label
$a_{d,j}$. The insertion comparison gives one row $u_c\in A_e$, chosen
before any spectator tuple or compact separated test $\theta$, whose
finite-cutoff tested values tend to the tested values $M_\theta$ of
$a_{d,j}$ on $S$. Write $P_{e,c}$ and $P_e$ for the finite and limiting
projections in the kernel construction, and $f_{\theta,c},f_\theta$ for
the corresponding tested linear functionals on $A_e$. Their convergence
and exact kernel identity are the hypotheses of
\eqref{eq:ope-test-kernel}. The subpower local equation input applies
also to
\[
 \frac{u_c-P_{e,c}u_c}{1+\|P_{e,c}u_c\|}\in\ker P_{e,c}.
\]
Indeed this is exactly the normalized residual treated following
\eqref{eq:ope-normalized-tests}: its first term has subpower bare norm
by the relay bound, and its second is bounded in reference coordinates,
hence has subpower bare norm by \eqref{eq:ope-bounded-bare}. Thus both
limits in \eqref{eq:ope-normalized-tests} hold, on the same filter and
with the same row $u_c$. The compactness argument there supplies one
$b_{de,j}\in\operatorname{ran}P_e$ with
$f_\theta(b_{de,j})=M_\theta$ for every $\theta$. This is
\eqref{eq:phase-inclusion-tested-input}. Its one compact refinement is
made for this row, not separately for each test. Applying it to the basis
labels supplies the raw maps $K_{de}$ above. This attachment uses only the
inclusion comparison; it asserts neither a derivative identity nor the
production of the retained analytic hypotheses.

\begin{proposition}[The same inclusions in every whole-word character]
\label{prop:phase-inclusion-every-character}
Let $S$, $A_d$ and the raw maps $K_{de}$ satisfy the hypotheses of
Proposition~\ref{prop:phase-field-inclusions}. Assume in addition the
unselected Hilbert space, bounded slab insertions, translation and
permutation identities, and placement-independent compressed operators
$M(w)$ of Section~\ref{ex-phases} for this same $S$. Let
$S_\chi(w)=\chi(M(w))$, $\chi\in\Xi$, be its whole-word character
tables as in \eqref{ex-phase-word}. Assume these tables are continuous
on each distinct-point configuration space, as supplied, with its own
hypotheses, by Proposition~\ref{ex-phase-caps}.
Then, for every character, every $d\le e$, every $a\in A_d$, every
finite raw spectator tuple $\boldsymbol b$ of arbitrary grades, and every
$x\in\operatorname{Conf}_{n+1}(\mathbb R^4)$,
\begin{equation}\label{eq:phase-inclusion-character}
 S_\chi(x;(e,K_{de}a),\boldsymbol b)
       =S_\chi(x;(d,a),\boldsymbol b).
\end{equation}
The identity holds in any slot and after integration against every smooth
compactly supported test in configuration space. The maps $K_{de}$ are
those fixed for $S$, independent of the character and of the exterior.
They induce coherent injective inclusions between the phase field
quotients and a common algebraic union in each phase. The natural maps
from the original quotients and union to these phase quotients and union
commute with all inclusions. If, for a given character,
$N_d=N_{\chi,d}$ for every $d$, these comparison maps are isomorphisms.
\end{proposition}
\begin{proof}
Proposition~\ref{prop:phase-field-inclusions} has already converted the
tested input on $S$ to the pointwise all-exterior identity
\eqref{eq:phase-inclusion-raw-response}. Fix $d,e,a,x,\boldsymbol b$,
and let $w$ and $w'$ be the words at these same positions, with distinguished
labels $(d,a)$ and $(e,K_{de}a)$ respectively. Translate both into the
same open slab of width $h>0$. This is possible because their common
finite set of times is bounded: choose the translation so that its
smallest time is positive and then choose $h$ larger than its largest
time. Denote the two bounded insertion operators by $B$ and $B'$.

Use the matrix-element identity from the proof of
Lemma~\ref{ex-phase-original-null}. For any positive-time words $u,z$,
the two quantities
\[
 \langle[u],B[z]\rangle,\qquad
 \langle[u],B'[z]\rangle
\]
are unselected correlations obtained by placing a reflected copy of $u$
below the slab and a shifted copy of $z$ above it. The exterior words
and all their labels are the same in both expressions. Their points are
internally distinct, and the three time ranges are disjoint; hence these
are legitimate separated exteriors. Translation and permutation
invariance of $S$ put the distinguished label first and undo the common
translation. Equation~\eqref{eq:phase-inclusion-raw-response}, with the
original spectators and these additional exterior labels, makes the two
matrix elements equal. This step uses all finite spectators, not merely
a finite family of two-point tests.

Sesquilinearity gives equality for finite linear combinations of
positive-time word vectors on both sides. Those vectors span a dense
subspace of the unselected Hilbert space. Since $B-B'$ is bounded,
continuity gives $\langle v,(B-B')z\rangle=0$ for all Hilbert vectors
$v,z$. Taking $v=(B-B')z$ proves $(B-B')z=0$ for every $z$, so $B=B'$.
Compression to the ground subspace preserves this equality. Placement
independence identifies the compressed operators with $M(w)$ and $M(w')$.
Applying any character proves \eqref{eq:phase-inclusion-character}.
There is no exceptional set of characters. Equivalently, one may apply a
continuous linear extension of the character after compression; only its
values on these whole-word operators are used. No multiplicativity for
separated subwords, positivity of such an extension, or differentiability
of a slab insertion is needed for this step.

Multilinearity and permutation symmetry of the slab operators, obtained
by the same dense matrix-element test from those properties of $S$,
pass to every $S_\chi$.
It moves the distinguished slot to any prescribed slot. Continuity on
configuration space makes each side integrable against a compactly
supported smooth test there; integrating the pointwise identity proves
the tested identity. Empty spectator tuples are included in this argument;
the comparison then concerns one insertion. The empty word has no slot
to raise and keeps its original normalization.

For clarity, this also gives the identity formulated using the original
quotient labels, without imposing linearity on a choice of raw
representatives. If $v=\pi_da$ and $r_d(v)$ is any representative of $v$,
then $r_d(v)-a\in N_d$. Similarly,
\[
 r_e(\iota_{de}v)-K_{de}a\in N_e,
 \qquad \iota_{de}v=\pi_eK_{de}a.
\]
Lemma~\ref{ex-phase-original-null} shows that both differences vanish in
every character correlation against arbitrary spectators. The same is
true for any changes of spectator representatives. Thus the identity on
the fixed original quotient plan and the raw identity
\eqref{eq:phase-inclusion-character} agree. The linear map is induced by
$K_{de}$ on the quotient, not by the representative section $r_d$.

Finally apply the construction of
Proposition~\ref{prop:phase-field-inclusions} to the phase response maps:
\eqref{eq:phase-inclusion-character} gives
$\mathcal R_{\chi,e}K_{de}=\mathcal R_{\chi,d}$ directly. The descent,
injectivity, uniqueness and composition proof there uses this identity,
linearity and symmetry only. Its earlier smoothness assumption was used
to obtain the pointwise identity from the tested one; that step is already
proved here. Proposition~\ref{prop:phase-field-union} therefore constructs
the phase union. By Lemma~\ref{ex-phase-original-null},
$N_d\subseteq N_{\chi,d}$, so the two-table comparison preceding this
subsection applies to these same raw maps. Explicitly, with
$q_{\chi,d}\pi_da=\pi_{\chi,d}a$,
\[
 q_{\chi,e}\iota_{de}\pi_da
 =\pi_{\chi,e}K_{de}a
 =\iota^\chi_{de}q_{\chi,d}\pi_da.
\]
Surjectivity of $\pi_d$ gives the commuting identity on $E_d$ and hence
on its algebraic union. If the null spaces are equal in every grade,
the $q_{\chi,d}$ and their compatible union map are isomorphisms as
proved there. This last conclusion compares field spaces and their
inclusions; it does not assert equality of numerical correlations of
$S$ and $S_\chi$.
\end{proof}

This supplies the same-map selected-table attachment conditional on the
stated unselected realization and whole-word construction. It does not
supply the local service, calibrated comparison, rank or ordered-limit
inputs in Paragraph~\ref{par:phase-inclusion-attachment}. Equality of the
null spaces, when needed for isomorphisms, retains the separate hypotheses
of the earlier phase-null argument. Derivative and Ward identities and
compatibility with additional field actions remain separate statements.
